% GENERATED FILE. Edit canonical lessons, then run npm run generate:books.
% canonical-source-hash: fnv1a64:e18b86293c8bb7e7
% canonical-lessons: LA-C398-violo, LA-C398-praesto, LA-C398-obligo, LA-C398-ratus, LA-C398-irritus

\chapter{Keeping My Word}
\label{ch:la-keeping-my-word}
\hlchaptermodality{398}

\begin{chapteropening}
\textbf{By the end of this chapter:} \emph{I can say that someone broke or kept his word and bound himself, and that an agreement holds good or comes to nothing.}

Complete the last lesson of chapter 398: I can say that someone broke or kept his word and bound himself, and that an agreement holds good or comes to nothing.
\end{chapteropening}

\section[violō, violāre]{violō, violāre — to violate, to break; to wrong}

\label{lesson:LA-C398-violo}



\begin{quote}
Before the new one: say the Latin for a pact, then the Latin for to bargain for.
\end{quote}

\subsection*{You'll want to know: violō, violāre}
\textbf{violō, violāre} — \textquotedblleft{}to violate, to break; to wrong\textquotedblright{}.

\begin{quote}
Hostēs foedus violāvērunt.
\end{quote}

\begin{quote}
\emph{The enemy broke the treaty.}
\end{quote}

\emph{violō, violāre, violāvī, violātum}, kin to \emph{vīs}, force. \emph{Fidem violāre}, to break one's word. English \emph{violate} and \emph{violence}.

\begin{cousinweb}
\emph{violō} is built on \emph{vīs} \textquotedblleft{}force\textquotedblright{} (EtymDB; Lewis \& Short): to treat with force. English \emph{violate} and \emph{violence} come from its family.
\end{cousinweb}

\begin{grammarlens}[title={What you've built}]
One more word for telling, explaining and describing at B1.
\end{grammarlens}

\subsection*{Your turn}
\begin{itemize}
\raggedright
  \item \emph{Say it:} \emph{violō, violāre}
  \item \emph{Say it:} \emph{violō, violāre}, once more
  \item \emph{Say it:} a sentence of your own with \emph{violō, violāre}
  \item \emph{Recall:} say the Latin for a pact, then the Latin for to bargain for, then say \emph{violō, violāre} again
\end{itemize}

\begin{tcolorbox}[breakable,colback=teal!4,colframe=teal!35!black,title={Before you move on}]
What does violō, violāre mean? (\textquotedblleft{}To violate, to break; to wrong\textquotedblright{}.) Say it once more.
\end{tcolorbox}
\section[praestō, praestāre]{praestō, praestāre — to make good, to fulfil; to excel}

\label{lesson:LA-C398-praesto}



\begin{quote}
Before the new one: say the Latin for to bargain for, then the Latin for to violate.
\end{quote}

\subsection*{You'll want to know: praestō, praestāre}
\textbf{praestō, praestāre} — \textquotedblleft{}to make good, to fulfil; to excel\textquotedblright{}.

\begin{quote}
Fidem praestitit.
\end{quote}

\begin{quote}
\emph{He kept his word.}
\end{quote}

\emph{praestō, praestāre, praestitī, praestitum}. \emph{Fidem praestāre}, to make one's word good; with the dative, to stand out above: \emph{omnibus praestat}.

\begin{cousinweb}
\emph{praestō} is built from \emph{prae} \textquotedblleft{}before, in front\textquotedblright{} and \emph{stō} \textquotedblleft{}to stand\textquotedblright{} (Lewis \& Short): to stand out in front, then to stand good for a promise.
\end{cousinweb}

\begin{grammarlens}[title={What you've built}]
One more word for telling, explaining and describing at B1.
\end{grammarlens}

\subsection*{Your turn}
\begin{itemize}
\raggedright
  \item \emph{Say it:} \emph{praestō, praestāre}
  \item \emph{Say it:} \emph{praestō, praestāre}, once more
  \item \emph{Say it:} a sentence of your own with \emph{praestō, praestāre}
  \item \emph{Recall:} say the Latin for to bargain for, then the Latin for to violate, then say \emph{praestō, praestāre} again
\end{itemize}

\begin{tcolorbox}[breakable,colback=teal!4,colframe=teal!35!black,title={Before you move on}]
What does praestō, praestāre mean? (\textquotedblleft{}To make good, to fulfil; to excel\textquotedblright{}.) Say it once more.
\end{tcolorbox}
\section[obligō, obligāre]{obligō, obligāre — to oblige, to bind}

\label{lesson:LA-C398-obligo}



\begin{quote}
Before the new one: say the Latin for to violate, then the Latin for to make good.
\end{quote}

\subsection*{You'll want to know: obligō, obligāre}
\textbf{obligō, obligāre} — \textquotedblleft{}to oblige, to bind\textquotedblright{}.

\begin{quote}
Prōmissō mē obligāvī.
\end{quote}

\begin{quote}
\emph{I bound myself by a promise.}
\end{quote}

\emph{obligō, obligāre, obligāvī, obligātum}, \emph{ob} + \emph{ligō}, to tie. English \emph{oblige} and \emph{obligation}.

\begin{cousinweb}
\emph{obligō} is built from \emph{ob} and \emph{ligō} \textquotedblleft{}to tie\textquotedblright{} (EtymDB; Lewis \& Short): to tie someone to a duty. English \emph{oblige} and \emph{obligation} come from it.
\end{cousinweb}

\begin{grammarlens}[title={What you've built}]
One more word for telling, explaining and describing at B1.
\end{grammarlens}

\subsection*{Your turn}
\begin{itemize}
\raggedright
  \item \emph{Say it:} \emph{obligō, obligāre}
  \item \emph{Say it:} \emph{obligō, obligāre}, once more
  \item \emph{Say it:} a sentence of your own with \emph{obligō, obligāre}
  \item \emph{Recall:} say the Latin for to violate, then the Latin for to make good, then say \emph{obligō, obligāre} again
\end{itemize}

\begin{tcolorbox}[breakable,colback=teal!4,colframe=teal!35!black,title={Before you move on}]
What does obligō, obligāre mean? (\textquotedblleft{}To oblige, to bind\textquotedblright{}.) Say it once more.
\end{tcolorbox}
\section[ratus, rata, ratum]{ratus, rata, ratum — fixed, settled, valid}

\label{lesson:LA-C398-ratus}



\begin{quote}
Before the new one: say the Latin for to make good, then the Latin for to oblige.
\end{quote}

\subsection*{You'll want to know: ratus, rata, ratum}
\textbf{ratus, rata, ratum} — \textquotedblleft{}fixed, settled, valid\textquotedblright{}.

\begin{quote}
Pactum ratum est.
\end{quote}

\begin{quote}
\emph{The pact is valid.}
\end{quote}

\emph{ratus, rata, ratum}, the participle of \emph{reor}, to reckon. \emph{Ratum habēre}, to ratify; \emph{prō ratā parte}, in proportion, gives English \emph{pro rata}.

\begin{cousinweb}
\emph{ratus} is the participle of \emph{reor} \textquotedblleft{}to reckon, to think\textquotedblright{} (EtymDB; Lewis \& Short): reckoned, so settled. English \emph{pro rata} and \emph{rate} come from it.
\end{cousinweb}

\begin{grammarlens}[title={What you've built}]
One more word for telling, explaining and describing at B1.
\end{grammarlens}

\subsection*{Your turn}
\begin{itemize}
\raggedright
  \item \emph{Say it:} \emph{ratus, rata, ratum}
  \item \emph{Say it:} \emph{ratus, rata, ratum}, once more
  \item \emph{Say it:} a sentence of your own with \emph{ratus, rata, ratum}
  \item \emph{Recall:} say the Latin for to make good, then the Latin for to oblige, then say \emph{ratus, rata, ratum} again
\end{itemize}

\begin{tcolorbox}[breakable,colback=teal!4,colframe=teal!35!black,title={Before you move on}]
What does ratus, rata, ratum mean? (\textquotedblleft{}Fixed, settled, valid\textquotedblright{}.) Say it once more.
\end{tcolorbox}
\section[irritus, irrita, irritum]{irritus, irrita, irritum — void, invalid; in vain}

\label{lesson:LA-C398-irritus}



\begin{quote}
Before the new one: say the Latin for to oblige, then the Latin for fixed.
\end{quote}

\subsection*{You'll want to know: irritus, irrita, irritum}
\textbf{irritus, irrita, irritum} — \textquotedblleft{}void, invalid; in vain\textquotedblright{}.

\begin{quote}
Cōnsilia eius irrita fuērunt.
\end{quote}

\begin{quote}
\emph{His plans came to nothing.}
\end{quote}

\emph{irritus, irrita, irritum} (also \emph{inritus}): \emph{in-} not + \emph{ratus}, fixed. \emph{Irritum fierī}, to come to nothing.

\begin{cousinweb}
\emph{irritus} is built from \emph{in-} \textquotedblleft{}not\textquotedblright{} and \emph{ratus} \textquotedblleft{}fixed, valid\textquotedblright{} (EtymDB; Lewis \& Short): not settled, so void.
\end{cousinweb}

\begin{grammarlens}[title={What you've built}]
One more word for telling, explaining and describing at B1.
\end{grammarlens}

\subsection*{Your turn}
\begin{itemize}
\raggedright
  \item \emph{Say it:} \emph{irritus, irrita, irritum}
  \item \emph{Say it:} \emph{irritus, irrita, irritum}, once more
  \item \emph{Say it:} a sentence of your own with \emph{irritus, irrita, irritum}
  \item \emph{Recall:} say the Latin for to oblige, then the Latin for fixed, then say \emph{irritus, irrita, irritum} again
\end{itemize}

\begin{tcolorbox}[breakable,colback=teal!4,colframe=teal!35!black,title={Before you move on}]
What does irritus, irrita, irritum mean? (\textquotedblleft{}Void, invalid; in vain\textquotedblright{}.) Say it once more.
\end{tcolorbox}
